\newpage
\section{Limitations and next steps}
\begin{enumerate}
\item Prefix Backward Walk is a revised method definition. Its proposed layout and arithmetic examples are not implementation or benchmark results.
\item Candidate-target stopping is approximate. A large returned set can still omit the most relevant reference documents.
\item Balanced independent bits give useful count examples, but real embedding prefixes can be uneven and correlated, especially inside a cluster.
\item Matryoshka motivates testing prefixes; it does not guarantee that every shorter sign prefix preserves recall.
\item Sorted-key lookup avoids rescoring old rows, but construction and updates still have costs. No insertion or speed advantage over HNSW has been measured.
\end{enumerate}

\section{Conclusion}
\begin{enumerate}
\item The common observation is unchanged: query signs define an ideal binary pattern, and mismatches have costs determined by the query values.
\item Bitplanes uses those values to narrow a set of documents while keeping alternative branches.
\item Backward Walk starts from the query's prefix, removes its last bit when more candidates are needed, and scores only newly exposed rows.
\item The earlier key-enumeration experiment tested a different method. A new fixed-data comparison is required before saying where prefix Backward Walk wins or loses.
\end{enumerate}

\begin{thebibliography}{9}\small\raggedright
\bibitem{exa} Exa Team. \emph{How we built a web-scale vector database}. 2024.
\url{https://exa.ai/blog/building-web-scale-vector-db}.
\bibitem{hnsw} Y. Malkov and D. Yashunin. \emph{Efficient and robust approximate nearest neighbor search using Hierarchical Navigable Small World graphs}. 2016; revised 2018.
\url{https://arxiv.org/abs/1603.09320}.
\bibitem{mrl} A. Kusupati et al. \emph{Matryoshka Representation Learning}. 2022.
\url{https://arxiv.org/abs/2205.13147}.
\bibitem{nomic} Nomic AI. \emph{nomic-embed-text-v1.5}. Model documentation.
\url{https://huggingface.co/nomic-ai/nomic-embed-text-v1.5}.
\bibitem{msmarco} Microsoft. \emph{MS MARCO}. Passage collection and benchmark resources.
\url{https://microsoft.github.io/msmarco/}.
\end{thebibliography}
